%%%%%%%%%%%%%%%%%%%%%%%%%%%%%%%%%%%%%%%%%%%%%%%%%%%%%%%%%%%%%%%%%%%%
%% I, the copyright holder of this work, release this work into the
%% public domain. This applies worldwide. In some countries this may
%% not be legally possible; if so: I grant anyone the right to use
%% this work for any purpose, without any conditions, unless such
%% conditions are required by law.
%%%%%%%%%%%%%%%%%%%%%%%%%%%%%%%%%%%%%%%%%%%%%%%%%%%%%%%%%%%%%%%%%%%%

\documentclass{beamer}
\usetheme[faculty=ped,basePath=../fibeamer,logo=../fibeamer/logo/mu/Logo-UC-3.png]{fibeamer}
\graphicspath{{../resources/}}

\usepackage{algorithm}
\usepackage{algpseudocode}

\usepackage[utf8]{inputenc}
\usepackage[
  main=english, %% By using `czech` or `slovak` as the main locale
                %% instead of `english`, you can typeset the
                %% presentation in either Czech or Slovak,
                %% respectively.
  czech, slovak %% The additional keys allow foreign texts to be
]{babel}        %% typeset as follows:
%%
%%   \begin{otherlanguage}{czech}   ... \end{otherlanguage}
%%   \begin{otherlanguage}{slovak}  ... \end{otherlanguage}
%%
%% These macros specify information about the presentation
\title{Casting, input() y print()} %% that will be typeset on the
\subtitle{Semestre I-2026} %% title page.
\author{Andreina Cota Vidaurre}
%% These additional packages are used within the document:
\usepackage{ragged2e}  % `\justifying` text
\usepackage{booktabs}  % Tables
\usepackage{tabularx}
\usepackage{tikz}      % Diagrams
\usetikzlibrary{calc, shapes, backgrounds}
\usetikzlibrary{positioning, arrows.meta, shadows.blur, shapes.geometric}
\usepackage{amsmath, amssymb}
\usepackage{url}       % `\url`s
\usepackage{listings}  % Code listings
\usepackage{graphicx}
\usepackage{amsmath}
\usepackage[T1]{fontenc}
\usepackage{xcolor}
\usepackage{minted}
% \usepackage{adjustbox}


% Definicion de pseudocodigo en espaniol
\lstdefinelanguage{PseudocodigoES}{
    morekeywords={INICIO,FIN,PARA,HACER,SI,ENTONCES,FIN_SI,FIN_PARA,IMPRIMIR},
    sensitive=true,
    morecomment=[l]{//},
    morestring=[b]"
}

% Definicion de pseudocodigo en ingles
\lstdefinelanguage{PseudocodeEN}{
    morekeywords={BEGIN,END,FOR,DO,IF,THEN,END_IF,END_FOR,PRINT},
    sensitive=true,
    morecomment=[l]{//},
    morestring=[b]"
}

\lstset{
    basicstyle=\ttfamily\small,
    keywordstyle=\bfseries,
    columns=fullflexible,
    keepspaces=true,
    showstringspaces=false,
    frame=single,
    xleftmargin=0em,
    framexleftmargin=0em,
    framesep=2pt,
    gobble=4,
    resetmargins=true,
    aboveskip=0pt,
    belowskip=0pt
}

\lstdefinestyle{pythonstyle}{
    language=Python,
    basicstyle=\ttfamily\small,
    keywordstyle=\bfseries,
    commentstyle=\itshape,
    stringstyle=\ttfamily,
    showstringspaces=false,
    columns=fullflexible,
    keepspaces=true,
    frame=single,
    breaklines=true,
    xleftmargin=0em,
    framexleftmargin=0em,
    framesep=2pt,
    aboveskip=0pt,
    belowskip=0pt,
    gobble=4,
    resetmargins=true,
    emph={print,input,int,float,str,bool,len,range},
    emphstyle=\bfseries
}

\lstset{style=pythonstyle}

% \definecolor{green_python}{HTML}{52A736}

% \newcolumntype{Y}{>{\raggedright\arraybackslash}X}

\frenchspacing
\begin{document}
  \shorthandoff{-}
  \frame[c]{\maketitle}


    \begin{frame}[fragile,t]{Casting, \mintinline{python}{input()} y \mintinline{python}{print()}}
        \begin{itemize}
            \item Casting es la conversión de un dato de un tipo a otro
            \item Se usa cuando un dato necesita cambiar de forma para poder operar correctamente
            \item En Python se hace con funciones como:
        \end{itemize}
        \begin{minted}{python}
            int()
            float()
            str()
            bool()
        \end{minted}
    \end{frame}

    \begin{frame}[fragile,t]{Casting}
        \begin{minted}[xleftmargin=-2cm, frame=leftline, fontsize=\small]{python}
            numero = "25"
            numero = int(numero)
            print(numero + 5)
        \end{minted}
    \end{frame}

    \begin{frame}{\mintinline{python}{int()}}
         \begin{itemize}
            \item Convierte un valor a número entero
            \item Elimina la parte decimal si recibe un número decimal (float)
            \item Puede convertir texto si el texto representa un número entero válido
         \end{itemize}
    \end{frame}

    \begin{frame}[fragile,t]{\mintinline{python}{int()}}
         \begin{minted}[xleftmargin=-2cm, frame=leftline, fontsize=\small]{python}
            int("8")      # 8
            int(4.9)      # 4
            int(True)     # 1
            int(False)    # 0
        \end{minted}
        Errores comunes:
        \begin{minted}[xleftmargin=-2cm, frame=leftline, fontsize=\small]{python}
            int("3.5")    # error
            int("hola")   # error
        \end{minted}
    \end{frame}

    \begin{frame}{\mintinline{python}{float()}}
        \begin{itemize}
            \item Convierte un valor a número decimal
            \item Puede convertir números enteros y textos numéricos
        \end{itemize}
    \end{frame}

    \begin{frame}[fragile,t]{\mintinline{python}{float()}}
        \begin{minted}[xleftmargin=-2cm, frame=leftline, fontsize=\small]{python}
            float("8.2")   # 8.2
            float("5")     # 5.0
            float(7)       # 7.0
        \end{minted}
        Error común:
        \begin{minted}[xleftmargin=-2cm, frame=leftline, fontsize=\small]{python}
            float("tres")  # error
        \end{minted}
    \end{frame}

    \begin{frame}{\mintinline{python}{str()}}
        \begin{itemize}
            \item Convierte un dato a texto
            \item Es útil para mostrar resultados en pantalla
            \item Permite unir números con texto
        \end{itemize}
    \end{frame}

    \begin{frame}[fragile,t]{\mintinline{python}{str()}}
        \begin{minted}[xleftmargin=-2cm, frame=leftline, fontsize=\small]{python}
            str(25)        # "25"
            str(3.14)      # "3.14"
            str(True)      # "True"
        \end{minted}
        Ejemplo práctico
        \begin{minted}[xleftmargin=-2cm, frame=leftline, fontsize=\small]{python}
            edad = 20
            print("Su edad es " + str(edad))
        \end{minted}
    \end{frame}

    \begin{frame}[fragile,t]{\mintinline{python}{bool()}}
        \begin{itemize}
            \item Convierte un valor a \mintinline{python}{True o False}
            \item Muchos valores se interpretan como verdaderos
            \item Algunos valores vacíos o cero se interpretan como falsos (False)
        \end{itemize}
    \end{frame}

    \begin{frame}[fragile,t]{\mintinline{python}{bool()}}
        \begin{minted}[xleftmargin=-2cm, frame=leftline, fontsize=\small]{python}
            bool(1)        # True
            bool(0)        # False
            bool(3.5)      # True
            bool(0.0)      # False
            bool("hola")   # True
            bool("")       # False
        \end{minted}
    \end{frame}

    \begin{frame}[fragile,t]{Errores frecuentes}
        \begin{itemize}
            \item Creer que \mintinline{python}{input()} devuelve números
            \item Usar \mintinline{python}{int("3.7")}
            \item Pensar que \mintinline{python}{bool("False")} da \mintinline{python}{False}
            \item Confundir \mintinline{python}{str(5)} con el número $5$
        \end{itemize}
    \end{frame}

    \begin{frame}[fragile,t]{\mintinline{python}{bool()}}
        \begin{itemize}
            \item \mintinline{python}{0, 0.0, ""} suelen dar \mintinline{python}{False}
            \item Casi todo lo demás da \mintinline{python}{True}
        \end{itemize}
        \begin{minted}[xleftmargin=-2cm, frame=leftline, fontsize=\small]{python}
            bool("False")
        \end{minted}
        Esto da:
        \begin{minted}[xleftmargin=-2cm, frame=leftline, fontsize=\small]{python}
            True
        \end{minted}
        porque \mintinline{python}{"False"} es un texto no vacío
    \end{frame}

    \begin{frame}[fragile,t]{\mintinline{python}{bool()}}
        \begin{minted}[xleftmargin=-2cm, frame=leftline, fontsize=\small]{python}
            texto = input("Ingrese True o False: ")
            valor = texto == "True"
            print(valor)
        \end{minted}
        Otra opción:
        \begin{minted}[xleftmargin=-2cm, frame=leftline, fontsize=\small]{python}
            texto = input("Ingrese verdadero o falso: ")
            valor = texto == "verdadero"
            print(valor)
        \end{minted}
    \end{frame}

    % \begin{frame}[fragile,t]{\mintinline{python}{bool()}}
    %     \begin{minted}[xleftmargin=-2cm, frame=leftline, fontsize=\small]{python}
    %         texto = input("Ingrese si/no: ")

    %         if texto == "si":
    %             valor = True
    %         elif texto == "no":
    %             valor = False
    %         else:
    %             print("Entrada no valida")
    %     \end{minted}
    % \end{frame}

    \begin{frame}[fragile,t]{\mintinline{python}{input()} y \mintinline{python}{print()}}
        \begin{itemize}
            \item \mintinline{python}{input()} \textbf{recibe} datos
            \item \mintinline{python}{print()} \textbf{muestra} datos
        \end{itemize}
    \end{frame}

    
    \begin{frame}[fragile,t]{\mintinline{python}{input()}}
        \begin{itemize}
            \item Permite ingresar datos desde el teclado
            \item Se usa para que el usuario escriba el dato requerido
            \item El valor ingresado se guarda normalmente en una variable
            \item Lo que devuelve es \textbf{texto}
        \end{itemize}
    \end{frame}

    \begin{frame}[fragile,t]{\mintinline{python}{input()}}
        \begin{minted}[xleftmargin=-2cm, frame=leftline, fontsize=\small]{python}
            input("Ingrese su nombre: ")
        \end{minted}

        \begin{minted}[xleftmargin=-2cm, frame=leftline, fontsize=\small]{python}
            nombre = input("Ingrese su nombre: ")
        \end{minted}

        \begin{minted}[xleftmargin=-2cm, frame=leftline, fontsize=\small]{python}
            nombre = input()
        \end{minted}

        \begin{minted}[xleftmargin=-2cm, frame=leftline, fontsize=\small]{python}
            valor1 = input("Ingrese el valor: ")
            valor1 = int(valor1)
        \end{minted}
    \end{frame}

    \begin{frame}[fragile,t]{\mintinline{python}{print()}}
        \begin{itemize}
            \item Muestra información en pantalla
            \item Se usa para comunicar resultados, mensajes o valores
            \item Puede mostrar texto, números o variables
        \end{itemize}
    \end{frame}

    \begin{frame}[fragile,t]{\mintinline{python}{print()}}
        \begin{minted}[xleftmargin=-2cm, frame=leftline, fontsize=\small]{python}
            print("Bienvenido al sistema")
        \end{minted}
        \begin{minted}[xleftmargin=-2cm, frame=leftline, fontsize=\small]{python}
            print(5.6)
        \end{minted}
        \begin{minted}[xleftmargin=-2cm, frame=leftline, fontsize=\small]{python}
            valor = 34
            print(valor)
        \end{minted}
        \begin{minted}[xleftmargin=-2cm, frame=leftline, fontsize=\small]{python}
            print(True)
        \end{minted}
        \begin{minted}[xleftmargin=-2cm, frame=leftline, fontsize=\small]{python}
            valor_booleano = False
            print(valor_booleano)
        \end{minted}
    \end{frame}

    \begin{frame}[fragile,t]{\mintinline{python}{print()}}
        \begin{minted}[xleftmargin=-2cm, frame=leftline, fontsize=\small]{python}
            nombre = input("Como te llamas? ")
            print("Hola")
            print(nombre)
        \end{minted}

        \begin{itemize}
            \item \mintinline{python}{input("Como te llamas? ")} pide un dato al usuario
            \item el dato se guarda en la variable \mintinline{python}{nombre}
            \item \mintinline{python}{print("Hola")} muestra un mensaje
            \item \mintinline{python}{print(nombre)} muestra un valor almacenado
        \end{itemize}
    \end{frame}

    \begin{frame}[fragile,t]{Diferencias entre \mintinline{python}{input() y print()}}
        \begin{itemize}
            \item \mintinline{python}{input()} espera al usuario. El programa se detiene hasta que el usuario escriba algo
            \item \mintinline{python}{print()} muestra información y el programa continua
            \item \mintinline{python}{input()} normalmente se asigna a una variable
            \item \mintinline{python}{print()} se usa para mostrar variables o resultados
        \end{itemize}
    \end{frame}

    \begin{frame}[fragile,t]{\mintinline{python}{input()}}
        \begin{minted}[xleftmargin=-2cm, frame=leftline, fontsize=\small]{python}
            edad = input("Ingrese su edad: ")
            print(edad + 1)
        \end{minted}
        El código produce un error porque \mintinline{python}{input()} devuelve texto. 
        Lo correcto sería convertir el texto a entero:

        \begin{minted}[xleftmargin=-2cm, frame=leftline, fontsize=\small]{python}
            edad = int(input("Ingrese su edad: "))
            print(edad + 1)
        \end{minted}
    \end{frame}
    
\end{document}
